% !PRINTED 89
\weberpage{89}
%<<CONT
choose. This is a single-valued function $F(z)$ of the variable
\[
z = e^{2\pi iv},
\]
and if, as before, $q = e^{\pi i\omega}$, then it has the periodicity property:
\weq{2}{F(q^2z) = e^{-2\pi ia}F(z),}
from which, for every integer $\nu$, it follows that:
\[
F(q^{2\nu}z) = e^{-2\pi ia\nu}F(z).
\]
The function $F(z)$ becomes infinite for a finite $z$ only
when one of the factors $q^{2\nu}z - 1$ vanishes, and in
particular
\weq{3}{(z-1)F(z) = 1 \quad (\text{for } z = 1).}
But by (2)
\[
(q^{2\nu}z - 1)F(q^{2\nu}z) = e^{-2\pi ia\nu}F(z)(q^{2\nu}z - 1)
\]
and consequently, by (3),
\weq{4}{F(z)(q^{2\nu}z - 1) = e^{2\pi ia\nu} \quad (\text{for } z = q^{-2\nu}).}
If we set, under the assumption of convergence that is presently
 to be examined more closely,
\weq{5}{S(z) = \sum_{-\infty,\infty}^{\nu} \frac{e^{2\pi ia\nu}}{q^{2\nu}z - 1},}
then, replacing $z$ by $q^2z$ and $\nu$ by $\nu - 1$, we obtain:
\weq{6}{S(q^2z) = e^{-2\pi ia}S(z),}
and the difference $F(z) - S(z)$ would therefore be, considered as a function of $v$,
by (3) and (4) an integral doubly periodic function
 of the second kind. But by § 16, 6. such a function must be an exponential
 function, and hence
\[
F(z) = S(z) + Ce^{-2\pi imv},
\]
where $m$ is an integer and $C$ a constant. But from (2) and
(6), unless $C = 0$, it follows that:
\[
e^{-2\pi im\omega} = e^{-2\pi ia};
\]
hence $a = m\omega + n$, where $m, n$ are integers. But if $a$
is a period, then $F(z)$ reduces to a constant or an
exponential function. Otherwise $C = 0$ must hold, and there follows
the development:
\weq{7}{\frac{\vartheta'_{11}\vartheta_{11}(v+a)}{2i\pi\vartheta_{11}(v)\vartheta_{11}(a)} = \sum_{-\infty,+\infty}^{\nu} \frac{e^{2\pi i\nu a}}{q^{2\nu}e^{2\pi iv} - 1}.}
